\section{Summary of Contributions}
\label{sec:conclusion_summary}

This section revisits the four research questions introduced in Section~\ref{sec:research_questions}.
For each of them, it summarizes the answer provided by this manuscript.

Research Question~\ref{rq:latent_space} concerned the approaches proposed in the literature to analyze the internal workings of \glspl{nn}.
To answer it, we first made explicit the working hypothesis shared by most of this literature, the \ConceptHypothesisName.
According to this hypothesis, \glspl{nn} generalize by organizing their \gls{ls} around concepts.
We then organized existing methods into four families: probing, interpretable directions, \glspl{sae}, and clustering.
Each family studies the \gls{ls} from a different angle, but they share the goal of understanding how \glspl{lr} are structured.

Research Question~\ref{rq:taxonomy} concerned the organization of existing \gls{xrl} methods \cite{alaarabiou:hal-04685607, alaarabiou2026taxonomy}.
To answer it, we first clarified what an explanation is by placing the observer at the center of the explanatory process.
Building on this frame, we identified the subjects and sources of explanation specific to \gls{rl}.
The subject of explanation is the policy of the agent, which can be explained either at the level of a single decision or at the level of its strategy over time.
The sources of explanation are the policies, the trajectories, and the value functions.
From these elements, we proposed a taxonomy that accounts for the specificities of \gls{rl} and used it to review the state of the art.

Research Question~\ref{rq:mechanisms} concerned the use of \glspl{lr} to identify the mechanisms underlying the decisions of \gls{rl} agents.
To answer it, we introduced \gls{car}, an \gls{xrl} method based on the principles of \gls{me} \cite{alaarabiou2026car}.
\gls{car} trains several surrogate policies independently to reproduce the same agent.
It then clusters their \glspl{lr} and compares the resulting partitions.
This procedure relies on the \RepresentationConvergenceHypothesisName, according to which independently trained networks develop similar concepts for a given task.
Applied to six Atari environments, \gls{car} recovers similar clustering structures across surrogate policies.
The qualitative analysis shows that these clusters capture interpretable aspects of the task, such as episode progression, specific game situations, and behavioral patterns.

Research Question~\ref{rq:evaluation} concerned the objective quantification of explanation quality.
To answer it, we introduced three metrics.
The \gls{cu} measures the consistency of the clusters across independently trained surrogate policies.
The \gls{as} measures how far these clusters depart from the structure of the observations, of the actions, and of an untrained network.
The experiments show that these two properties evolve in opposite directions across the network.
The earliest layers are the most consistent, but also the least abstract.
The \gls{acu} combines both properties into a single measure and provides a practical criterion for selecting the layer to analyze.
These metrics make it possible to compare competing hypotheses about the structure of a \gls{ls} without relying on manual inspection.

Taken together, these contributions form a coherent chain.
The \MechanisticHypothesisName states that \glspl{nn} can be explained through their internal mechanisms.
The \ConceptHypothesisName states that these mechanisms are organized around concepts.
The \RepresentationConvergenceHypothesisName states that these concepts reflect the task.
\gls{car} translates these hypotheses into an empirical procedure, and the experiments provide support for them.
